\scanpage{32}
\begin{center}-0.2-\end{center}

projective $n$-space \uline{over $\mathbb Z$}, unless, in a particular chapter, all considerations are over a ground field $k$, in which case they will denote affine $n$-space and projective $n$-space over $k$. $\operatorname{PGL}(n)$ will be the projective group acting on $\mathbb P_n$, $\operatorname{GL}(n)$ the general linear group acting on $\mathbb A^n$, and $\mathbb G_m=\operatorname{GL}(1)$. The same conventions on the base scheme hold for these group schemes.

(6) $1_X$ is the identity morphism from $X$ to $X$. $p_1$ and $p_2$ are the projections from $X\times_S Y$ to $X$ and $Y$; $p_{12}$, etc. are the projections from $X\times_S Y\times_S Z$ to $X\times_S Y$, etc. If $f:X_1\longrightarrow X_2$ and $g:Y_1\longrightarrow Y_2$ are $S$-morphisms, then $f\times g:X_1\times_S Y_1\longrightarrow X_2\times_S Y_2$ is the product. If $f:X\longrightarrow Y_1$ and $g:X\longrightarrow Y_2$ are $S$-morphisms, then $(f,g):X\longrightarrow Y_1\times_S Y_2$ is the induced $X$-morphism.

(7) If $X$ is a pre-scheme, and $x\in X$, then $\unsure{\kappa|\mathcal K}(x)$ will denote the residue field of $\underline o_{x,X}$. If $X$ is a scheme over a field $k$, $X_k$ will denote the set of points $x\in X$ such that $k\xrightarrow{\sim}\unsure{\kappa|\mathcal K}(x)$.

\margin{\figure{\begin{tikzpicture}[scale=.5]\draw[gray!40] (0,0)--(.5,-.02) .. controls (1,-.05) and (1,.35) .. (.52,.4);\draw[gray!40] (.3,0)--(.31,.36);\end{tikzpicture}}}

\fbox{\S1.} \uline{Definition 0.1}: A group pre-scheme $G/S$ is a morphism $\pi:G\longrightarrow S$ of pre-schemes, plus $S$-morphisms $\mu:G\times_S G\longrightarrow G$, $\beta:G\longrightarrow G$, $e:S\longrightarrow G$ satisfying the usual identities:

(a) \uline{Associativity}:
\diagram{\begin{tikzcd}
G\times_S G\times_S G \arrow[r,"1_G\times\mu"] \arrow[d,"\mu\times1_G"'] & G\times_S G \arrow[d,"\mu"] \\
G\times_S G \arrow[r,"\mu"] & G
\end{tikzcd}} commutes.

\mspage{31}
